\documentclass[12pt]{article}
\usepackage[utf8]{inputenc}
\usepackage{amsmath, amssymb, amsthm}
\usepackage{geometry}
\usepackage{xcolor}
\usepackage[most]{tcolorbox}
\usepackage{enumitem}
\usepackage{fancyhdr}

% Page setup
\geometry{margin=1in}
\setlength{\headheight}{14.5pt}
\pagestyle{fancy}
\fancyhf{}
\rhead{AMC12/COMC Problem Analysis}
\lhead{\today}

% Color definitions
\definecolor{problemconfig}{RGB}{230, 242, 255}    % Light blue
\definecolor{strategic}{RGB}{255, 230, 242}       % Light pink
\definecolor{tactical}{RGB}{242, 255, 230}        % Light green
\definecolor{techniques}{RGB}{255, 242, 230}      % Light yellow
\definecolor{verification}{RGB}{255, 230, 230}    % Light red
\definecolor{solution}{RGB}{230, 255, 255}        % Light cyan
\definecolor{competition}{RGB}{242, 242, 255}     % Light purple
\definecolor{gold}{RGB}{218, 165, 32}

% Box styles
\newtcolorbox{problemconfigbox}{
    colback=problemconfig,
    colframe=blue,
    title=Problem Configuration Analysis,
    fonttitle=\bfseries,
    arc=3pt,
    boxrule=1pt,
    breakable
}

\newtcolorbox{strategicbox}{
    colback=strategic,
    colframe=purple,
    title=Strategic Framework Application,
    fonttitle=\bfseries,
    arc=3pt,
    boxrule=1pt,
    breakable
}

\newtcolorbox{tacticalbox}{
    colback=tactical,
    colframe=green,
    title=Tactical Methodology Selection,
    fonttitle=\bfseries,
    arc=3pt,
    boxrule=1pt,
    breakable
}

\newtcolorbox{techniquesbox}{
    colback=techniques,
    colframe=orange,
    title=Conversion Techniques Application,
    fonttitle=\bfseries,
    arc=3pt,
    boxrule=1pt,
    breakable
}

\newtcolorbox{verificationbox}{
    colback=verification,
    colframe=red,
    title=Verification Strategy Design,
    fonttitle=\bfseries,
    arc=3pt,
    boxrule=1pt,
    breakable
}

\newtcolorbox{solutionbox}{
    colback=solution,
    colframe=cyan,
    title=Solution Roadmap,
    fonttitle=\bfseries,
    arc=3pt,
    boxrule=1pt,
    breakable
}

\newtcolorbox{competitionbox}{
    colback=competition,
    colframe=violet,
    title=Competition Strategy Integration,
    fonttitle=\bfseries,
    arc=3pt,
    boxrule=1pt,
    breakable
}

\newtcolorbox{keyinsight}{
    colback=yellow!10,
    colframe=gold,
    title=\textcolor{gold}{$\star$ Key Insight},
    fonttitle=\bfseries,
    arc=3pt,
    boxrule=2pt,
    breakable
}

\newtcolorbox{warning}{
    colback=red!5,
    colframe=red!70!black,
    title=\textcolor{red!70!black}{$\triangle$ Warning},
    fonttitle=\bfseries,
    arc=3pt,
    boxrule=2pt,
    breakable
}

\newtcolorbox{protip}{
    colback=green!5,
    colframe=green!70!black,
    title=\textcolor{green!70!black}{$\checkmark$ Pro Tip},
    fonttitle=\bfseries,
    arc=3pt,
    boxrule=2pt,
    breakable
}

\begin{document}

\title{\textbf{Strategic Analysis: AMC 12B 2024 Problem 24}}
\author{Comprehensive Problem-Solving Framework}
\date{\today}
\maketitle

\section*{Problem Statement}

What is the number of ordered triples $(a,b,c)$ of positive integers, with $a \leq b \leq c \leq 9$, such that there exists a (non-degenerate) triangle $\triangle ABC$ with an integer inradius for which $a$, $b$, and $c$ are the lengths of the altitudes from $A$ to $\overline{BC}$, $B$ to $\overline{AC}$, and $C$ to $\overline{AB}$, respectively?

(Recall that the inradius of a triangle is the radius of the largest possible circle that can be inscribed in the triangle.)

$$\textbf{(A) } 2 \qquad \textbf{(B) } 3 \qquad \textbf{(C) } 4 \qquad \textbf{(D) } 5 \qquad \textbf{(E) } 6$$

\vspace{1cm}

\begin{problemconfigbox}
\subsection*{1. Problem Configuration Analysis}

\subsubsection*{A. Problem Classification}

\begin{itemize}[leftmargin=*]
    \item \textbf{Problem Type}: To Find (counting problem with constraints)
    \item \textbf{Difficulty Band}: Late-band (Problem 24/25)
    \item \textbf{Estimated Time}: 7-10 minutes
    \item \textbf{Competition Context}: AMC 12B, multiple choice format
    \item \textbf{Mathematical Domain}: Geometry (triangle properties, inradius) with Number Theory (integer constraints, Diophantine equations) and Combinatorics (counting)
    \item \textbf{Cross-Domain Nature}: This is a sophisticated cross-domain problem requiring:
    \begin{itemize}
        \item \textbf{Geometry}: Triangle altitudes, inradius, area formulas
        \item \textbf{Algebra}: Deriving relationship between inradius and altitudes
        \item \textbf{Number Theory}: Finding integer solutions to equations
        \item \textbf{Combinatorics}: Counting valid ordered triples
        \item \textbf{Inequalities}: Triangle inequality for reciprocals
    \end{itemize}
\end{itemize}

\subsubsection*{B. Problem Structure Identification}

\textbf{Given Information}:
\begin{itemize}
    \item Ordered triple $(a, b, c)$ of positive integers
    \item Constraint: $a \leq b \leq c \leq 9$
    \item $a, b, c$ are altitudes of a triangle $\triangle ABC$
    \item Triangle is non-degenerate
    \item Inradius $r$ must be a positive integer
    \item Multiple choice answers: 2, 3, 4, 5, 6
\end{itemize}

\textbf{Constraints}:
\begin{itemize}
    \item \textbf{Ordering}: $a \leq b \leq c$ (given)
    \item \textbf{Upper bound}: $c \leq 9$
    \item \textbf{Positive integers}: $a, b, c, r \in \mathbb{Z}^+$
    \item \textbf{Non-degenerate triangle}: Must satisfy triangle inequality
    \item \textbf{Altitude constraint}: For altitudes $a, b, c$ to form a valid triangle, the reciprocals $\frac{1}{a}, \frac{1}{b}, \frac{1}{c}$ must satisfy triangle inequality
\end{itemize}

\textbf{Objective}:
\begin{itemize}
    \item Count the number of valid ordered triples $(a, b, c)$
    \item Must derive relationship between altitudes and inradius
\end{itemize}

\textbf{Hidden Assumptions}:
\begin{itemize}
    \item Need to derive the formula relating $r$, $a$, $b$, $c$
    \item The triangle inequality applies to reciprocals of altitudes
    \item The number of solutions is small (answer choices 2-6)
    \item Systematic case analysis on $r$ is feasible
\end{itemize}

\textbf{Answer Choice Leverage}:
\begin{itemize}
    \item Small range (2-6) suggests few valid solutions
    \item Can use answer to guide search scope
    \item Odd answer (B) with value 3 stands out
    \item Suggests looking for pattern (possibly multiples)
\end{itemize}

\subsubsection*{C. Strategic Orientation}

\textbf{Hypothesis}: There exists a simple relationship between inradius $r$ and altitudes $a, b, c$ that can be exploited to find integer solutions.

\textbf{Conclusion}: Count all $(a, b, c)$ satisfying the derived constraint.

\textbf{Notation}:
\begin{itemize}
    \item $a, b, c$ = altitudes from vertices $A, B, C$ respectively
    \item $r$ = inradius of the triangle
    \item $s$ = semiperimeter of the triangle
    \item $K$ = area of the triangle
    \item Side lengths = $AB, BC, AC$ (to be determined from altitudes)
\end{itemize}

\textbf{Intuitive Approach}:
\begin{itemize}
    \item Derive formula: Use area and inradius formula
    \item Bound search: Use constraints to limit possible values of $r$
    \item Case analysis: Check each possible value of $r$ systematically
    \item Triangle inequality: Verify reciprocal condition for each candidate
\end{itemize}

\textbf{Cross-Domain Potential}:
\begin{itemize}
    \item Pure geometric derivation using area formulas
    \item Algebraic approach with reciprocal relationships
    \item Number theoretic search for integer solutions
    \item Computational enumeration with bounds
\end{itemize}

\end{problemconfigbox}

\begin{strategicbox}
\subsection*{2. Strategic Framework Application}

\subsubsection*{A. Psychological Strategies}

\textbf{Mental Toughness}:
\begin{itemize}
    \item This is Problem 24 - expect sophisticated reasoning
    \item The relationship between altitudes and inradius is not standard knowledge
    \item Must derive formula from first principles
    \item Be prepared for careful case analysis
\end{itemize}

\textbf{Creativity Cultivation}:
\begin{itemize}
    \item Think: How does inradius relate to area and perimeter?
    \item Explore: Can altitudes be expressed in terms of sides and area?
    \item Consider: What constraints does the triangle inequality impose?
    \item Question: What bounds can we place on $r$?
\end{itemize}

\textbf{Iterative Refinement}:
\begin{itemize}
    \item First: Derive the key formula relating $r$ to $a, b, c$
    \item Then: Establish bounds on $r$ from the constraints
    \item Next: Systematically check each case
    \item Finally: Verify triangle inequality for each candidate
\end{itemize}

\subsubsection*{B. Investigation Strategies}

\textbf{Startup Orientation}:
\begin{itemize}
    \item Recognize this as a counting problem with constraints
    \item Identify that a formula must be derived
    \item Note that answer choices suggest small count
\end{itemize}

\textbf{Get Your Hands Dirty}:
\begin{itemize}
    \item Try simple case: $(3, 3, 3)$ (equilateral in altitude space)
    \item Check: Does this give integer inradius?
    \item Explore: What about $(6, 6, 6)$ or $(9, 9, 9)$?
    \item Pattern: Multiples of 3 seem to work
\end{itemize}

\textbf{Penultimate Step}:
\begin{itemize}
    \item Final step: Count valid triples
    \item Penultimate step: Verify each candidate satisfies triangle inequality
    \item Working backward: What values of $r$ are possible given constraints?
\end{itemize}

\textbf{Wishful Thinking}:
\begin{itemize}
    \item Wish: "I want a simple formula like $\frac{1}{r} = \frac{1}{a} + \frac{1}{b} + \frac{1}{c}$"
    \item Reality: This is exactly the formula we'll derive!
    \item Wish: "I want only a few cases to check"
    \item Reality: Only $r = 1, 2, 3$ are possible
\end{itemize}

\textbf{Make It Easier}:
\begin{itemize}
    \item Focus on deriving the formula first
    \item Use bounds to limit case analysis
    \item Leverage symmetry (if $a = b = c$, calculations simplify)
\end{itemize}

\textbf{Conjecture via Small Cases}:
\begin{itemize}
    \item Test $(3, 3, 3)$: Does $\frac{1}{r} = \frac{3}{3} = 1$? Yes, $r = 1$ works!
    \item Test $(6, 6, 6)$: Does $\frac{1}{r} = \frac{3}{6} = \frac{1}{2}$? Yes, $r = 2$ works!
    \item Pattern: $(3k, 3k, 3k)$ gives $r = k$ for $k \leq 3$
\end{itemize}

\textbf{Visualization}:
\begin{itemize}
    \item Visualize triangle with altitudes
    \item Consider relationship between altitude and opposite side
    \item Think about how inradius scales with triangle size
\end{itemize}

\subsubsection*{C. Late-Band Strategic Playbook}

\textbf{Answer-Choice Leverage}:
\begin{itemize}
    \item Range 2-6 suggests exhaustive search is feasible
    \item Answer (B) = 3 suggests pattern with triples
    \item Can work backwards: if answer is 3, what are the triples?
\end{itemize}

\textbf{Make It Easier}:
\begin{itemize}
    \item Start with symmetric case $a = b = c$
    \item Use bounds to limit $r$ to small values
    \item Focus on deriving key formula first
\end{itemize}

\textbf{Penultimate Step}:
\begin{itemize}
    \item If we know $r \in \{1, 2, 3\}$, solve for each case
    \item Check which give valid $(a, b, c)$ with $c \leq 9$
\end{itemize}

\textbf{Wishful Thinking}:
\begin{itemize}
    \item Assume the pattern is simple (multiples of 3)
    \item Verify this guess systematically
\end{itemize}

\textbf{Conjecture via Pattern}:
\begin{itemize}
    \item Notice $(3, 3, 3)$, $(6, 6, 6)$, $(9, 9, 9)$ are candidates
    \item These are all the multiples of 3 up to 9
    \item Hypothesis: Answer is 3
\end{itemize}

\end{strategicbox}

\begin{tacticalbox}
\subsection*{3. Tactical Methodology Selection}

\subsubsection*{A. Core Mathematical Tactics}

\textbf{Primary Tactics}:
\begin{enumerate}
    \item \textbf{Inradius Formula}:
    $$r = \frac{K}{s}$$
    where $K$ is the area and $s$ is the semiperimeter
    
    \item \textbf{Area from Altitude}:
    If $h_a$ is the altitude to side $a'$, then:
    $$K = \frac{1}{2} a' \cdot h_a$$
    
    \item \textbf{Derivation of Key Formula}:
    Combining area and inradius relations:
    $$\frac{1}{r} = \frac{1}{a} + \frac{1}{b} + \frac{1}{c}$$
    
    \item \textbf{Triangle Inequality for Reciprocals}:
    For altitudes $a, b, c$ to form a valid triangle:
    $$\frac{1}{b} + \frac{1}{c} > \frac{1}{a}$$
    
    \item \textbf{Systematic Case Analysis}:
    Check each possible value of $r$ within bounds
    
    \item \textbf{Bound Analysis}:
    From $a \leq b \leq c \leq 9$:
    $$\frac{1}{r} \geq \frac{1}{9} + \frac{1}{9} + \frac{1}{9} = \frac{1}{3}$$
    Therefore $r \leq 3$
\end{enumerate}

\textbf{Supporting Tactics}:
\begin{itemize}
    \item \textbf{Integer Solutions}: Finding positive integer $(a, b, c, r)$ satisfying the equation
    \item \textbf{Ordered Triples}: Ensuring $a \leq b \leq c$
    \item \textbf{Feasibility Check}: Verifying triangle inequality
    \item \textbf{Exhaustive Enumeration}: Within bounded search space
\end{itemize}

\subsubsection*{B. Crossover Techniques}

\begin{itemize}
    \item \textbf{Geometry $\to$ Algebra}: Area formulas to algebraic equation
    \item \textbf{Number Theory}: Finding integer solutions to Diophantine-like equation
    \item \textbf{Inequalities}: Triangle inequality translated to reciprocals
    \item \textbf{Combinatorics}: Counting valid configurations
    \item \textbf{Bounds}: Using constraints to limit search space
\end{itemize}

\end{tacticalbox}

\begin{techniquesbox}
\subsection*{4. Conversion Techniques Application}

\subsubsection*{A. Recognition Index}

\textbf{Triggered Techniques}:
\begin{itemize}
    \item \textbf{Inradius-Area Formula}: $K = rs$ where $s$ is semiperimeter
    \item \textbf{Altitude-Area Formula}: $K = \frac{1}{2} \cdot \text{base} \cdot \text{height}$
    \item \textbf{Reciprocal Relationships}: Working with $\frac{1}{a}, \frac{1}{b}, \frac{1}{c}$
    \item \textbf{Harmonic Mean}: The formula has harmonic mean structure
    \item \textbf{Bounded Search}: Integer constraints limit possibilities
\end{itemize}

\subsubsection*{B. Technique Selection Priority}

\textbf{Primary Path} (Recommended):
\begin{enumerate}
    \item Derive the key formula: $\frac{1}{r} = \frac{1}{a} + \frac{1}{b} + \frac{1}{c}$
    \item Establish bounds: $r \leq 3$ from constraint $c \leq 9$
    \item Case analysis on $r = 1, 2, 3$
    \item For each $r$, find all $(a, b, c)$ with $a \leq b \leq c \leq 9$
    \item Verify triangle inequality: $\frac{1}{b} + \frac{1}{c} > \frac{1}{a}$
    \item Count valid triples
\end{enumerate}

\textbf{Alternative Path} (Pattern Recognition):
\begin{enumerate}
    \item Guess that symmetric cases $(a, a, a)$ work
    \item Test $(3, 3, 3)$, $(6, 6, 6)$, $(9, 9, 9)$
    \item Verify these give $r = 1, 2, 3$
    \item Check if any non-symmetric cases work
    \item Conclude answer is 3
\end{enumerate}

\subsubsection*{C. Integration Strategy}

\textbf{Conversion Sequence}:
$$\text{Triangle Properties} \xrightarrow{\text{Area Formulas}} \text{Inradius-Altitude Relation}$$
$$\xrightarrow{\text{Algebra}} \frac{1}{r} = \frac{1}{a} + \frac{1}{b} + \frac{1}{c} \xrightarrow{\text{Bounds}} r \leq 3$$
$$\xrightarrow{\text{Case Analysis}} \text{Integer Solutions} \xrightarrow{\text{Triangle Inequality}} \text{Valid Triples}$$

\textbf{Key Derivation}:

From $K = rs$ where $s = \frac{AB + BC + AC}{2}$:
$$\frac{r}{c} = \frac{AB}{AB + BC + AC}$$

Similarly:
$$\frac{r}{b} = \frac{AC}{AB + BC + AC}, \quad \frac{r}{a} = \frac{BC}{AB + BC + AC}$$

Adding:
$$\frac{r}{a} + \frac{r}{b} + \frac{r}{c} = \frac{AB + BC + AC}{AB + BC + AC} = 1$$

Therefore:
$$\boxed{\frac{1}{r} = \frac{1}{a} + \frac{1}{b} + \frac{1}{c}}$$

\subsubsection*{D. Conflict Resolution}

\textbf{Potential Conflicts}:
\begin{itemize}
    \item Multiple solutions for same $r$
    \item Triangle inequality violations
    \item Boundary cases at $c = 9$
\end{itemize}

\textbf{Resolution Strategy}:
\begin{itemize}
    \item Systematically check all divisors and factorizations
    \item Always verify triangle inequality
    \item Be careful with ordered constraint $a \leq b \leq c$
\end{itemize}

\end{techniquesbox}

\begin{verificationbox}
\subsection*{5. Verification Strategy Design}

\subsubsection*{A. Verification Level Selection}

\textbf{Recommended Level}: Level 4-5 (3-5 minutes)

\textbf{Rationale}:
\begin{itemize}
    \item Problem 24 is high-value (late-band)
    \item Derivation must be verified
    \item Each case needs checking
    \item Triangle inequality must be verified for each triple
    \item Small answer count makes exhaustive verification feasible
\end{itemize}

\subsubsection*{B. Specialized Verification Methods}

\textbf{For Each Candidate Triple $(a, b, c)$}:

\begin{enumerate}
    \item \textbf{Verify Formula}:
    $$\frac{1}{r} = \frac{1}{a} + \frac{1}{b} + \frac{1}{c}$$
    Check that this gives an integer $r$
    
    \item \textbf{Verify Bounds}:
    $$a \leq b \leq c \leq 9$$
    
    \item \textbf{Verify Triangle Inequality}:
    $$\frac{1}{b} + \frac{1}{c} > \frac{1}{a}$$
    This ensures a non-degenerate triangle exists
    
    \item \textbf{Verify Counting}:
    Ensure all cases are covered and no duplicates
\end{enumerate}

\textbf{Complete Verification for Found Solutions}:

\textbf{Triple 1: $(3, 3, 3)$}
\begin{align*}
\frac{1}{r} &= \frac{1}{3} + \frac{1}{3} + \frac{1}{3} = 1 \implies r = 1 \quad \checkmark \\
\text{Bounds:} &\quad 3 \leq 3 \leq 3 \leq 9 \quad \checkmark \\
\text{Triangle:} &\quad \frac{1}{3} + \frac{1}{3} = \frac{2}{3} > \frac{1}{3} \quad \checkmark
\end{align*}

\textbf{Triple 2: $(6, 6, 6)$}
\begin{align*}
\frac{1}{r} &= \frac{1}{6} + \frac{1}{6} + \frac{1}{6} = \frac{1}{2} \implies r = 2 \quad \checkmark \\
\text{Bounds:} &\quad 6 \leq 6 \leq 6 \leq 9 \quad \checkmark \\
\text{Triangle:} &\quad \frac{1}{6} + \frac{1}{6} = \frac{1}{3} > \frac{1}{6} \quad \checkmark
\end{align*}

\textbf{Triple 3: $(9, 9, 9)$}
\begin{align*}
\frac{1}{r} &= \frac{1}{9} + \frac{1}{9} + \frac{1}{9} = \frac{1}{3} \implies r = 3 \quad \checkmark \\
\text{Bounds:} &\quad 9 \leq 9 \leq 9 \leq 9 \quad \checkmark \\
\text{Triangle:} &\quad \frac{1}{9} + \frac{1}{9} = \frac{2}{9} > \frac{1}{9} \quad \checkmark
\end{align*}

\textbf{Verification of Exhaustiveness}:
\begin{itemize}
    \item Case $r = 1$: Only $(3, 3, 3)$ works (verified in solution)
    \item Case $r = 2$: Only $(6, 6, 6)$ works (verified in solution)
    \item Case $r = 3$: Only $(9, 9, 9)$ works (verified in solution)
    \item Case $r \geq 4$: Impossible since $\frac{1}{r} < \frac{1}{3} \leq \frac{1}{c} + \frac{1}{c} + \frac{1}{c}$
\end{itemize}

\subsubsection*{C. Confidence Calibration}

\textbf{Confidence Assessment}: 95-98\%

\textbf{Reasons for High Confidence}:
\begin{itemize}
    \item Formula derived from first principles
    \item Systematic case analysis leaves no gaps
    \item All three solutions verified independently
    \item Triangle inequality checked for each
    \item Bounds argument shows no other cases possible
    \item Answer matches choice (B)
\end{itemize}

\textbf{Residual Uncertainty}:
\begin{itemize}
    \item Possible arithmetic error in case analysis (2\%)
    \item Might have missed a non-symmetric case (3\%)
\end{itemize}

\end{verificationbox}

\begin{keyinsight}
\textbf{Central Insight}: The key relationship $\frac{1}{r} = \frac{1}{a} + \frac{1}{b} + \frac{1}{c}$ connects the inradius to the altitudes through a harmonic-like sum. Combined with the bounds $a \leq b \leq c \leq 9$, this limits $r$ to at most 3. Systematic case analysis reveals that only the symmetric triples $(3, 3, 3)$, $(6, 6, 6)$, and $(9, 9, 9)$ satisfy all constraints, giving exactly 3 valid solutions.
\end{keyinsight}

\begin{solutionbox}
\subsection*{6. Solution Roadmap}

\subsubsection*{Step-by-Step Solution}

\textbf{Step 1: Derive the Key Formula}

We need to relate the inradius $r$ to the altitudes $a$, $b$, $c$.

Recall: For a triangle with area $K$ and semiperimeter $s$:
$$K = rs$$

For a triangle with sides $AB$, $BC$, $AC$ and altitudes $c$, $a$, $b$ respectively:
$$K = \frac{1}{2} \cdot AB \cdot c = \frac{1}{2} \cdot BC \cdot a = \frac{1}{2} \cdot AC \cdot b$$

Therefore:
$$AB = \frac{2K}{c}, \quad BC = \frac{2K}{a}, \quad AC = \frac{2K}{b}$$

The semiperimeter is:
$$s = \frac{AB + BC + AC}{2} = \frac{2K}{c} + \frac{2K}{a} + \frac{2K}{b} \cdot \frac{1}{2} = K\left(\frac{1}{a} + \frac{1}{b} + \frac{1}{c}\right)$$

Wait, let me recalculate:
$$s = \frac{AB + BC + AC}{2} = \frac{1}{2}\left(\frac{2K}{c} + \frac{2K}{a} + \frac{2K}{b}\right) = K\left(\frac{1}{c} + \frac{1}{a} + \frac{1}{b}\right)$$

From $K = rs$:
$$K = r \cdot K\left(\frac{1}{a} + \frac{1}{b} + \frac{1}{c}\right)$$

Dividing by $K$ (assuming $K \neq 0$):
$$1 = r\left(\frac{1}{a} + \frac{1}{b} + \frac{1}{c}\right)$$

Therefore:
$$\boxed{\frac{1}{r} = \frac{1}{a} + \frac{1}{b} + \frac{1}{c}}$$

\textbf{Step 2: Establish Bounds on $r$}

Given $a \leq b \leq c \leq 9$, we have:
$$\frac{1}{r} = \frac{1}{a} + \frac{1}{b} + \frac{1}{c} \leq \frac{1}{a} + \frac{1}{a} + \frac{1}{a} = \frac{3}{a}$$

Since $a \geq 1$:
$$\frac{1}{r} \leq 3 \implies r \text{ can be arbitrarily large}$$

This doesn't help. Let's use the other direction:
$$\frac{1}{r} = \frac{1}{a} + \frac{1}{b} + \frac{1}{c} \geq \frac{1}{c} + \frac{1}{c} + \frac{1}{c} = \frac{3}{c}$$

Since $c \leq 9$:
$$\frac{1}{r} \geq \frac{3}{9} = \frac{1}{3}$$

Therefore:
$$r \leq 3$$

Since $r$ must be a positive integer: $\boxed{r \in \{1, 2, 3\}}$

\textbf{Step 3: Triangle Inequality Condition}

For a non-degenerate triangle to exist with altitudes $a$, $b$, $c$, the reciprocals must satisfy the triangle inequality:
$$\frac{1}{a} + \frac{1}{b} > \frac{1}{c}, \quad \frac{1}{b} + \frac{1}{c} > \frac{1}{a}, \quad \frac{1}{c} + \frac{1}{a} > \frac{1}{b}$$

Since $a \leq b \leq c$, we have $\frac{1}{a} \geq \frac{1}{b} \geq \frac{1}{c}$. 

The critical constraint is:
$$\boxed{\frac{1}{b} + \frac{1}{c} > \frac{1}{a}}$$

The other two inequalities are automatically satisfied when this one is.

\textbf{Step 4: Case Analysis}

\textbf{Case 1: $r = 1$}

We need: $\frac{1}{a} + \frac{1}{b} + \frac{1}{c} = 1$ with $a \leq b \leq c \leq 9$

Since $\frac{1}{a} \geq \frac{1}{b} \geq \frac{1}{c}$:
$$1 = \frac{1}{a} + \frac{1}{b} + \frac{1}{c} \leq \frac{3}{a}$$

Therefore $a \leq 3$.

\textbf{Subcase $a = 1$}: $\frac{1}{b} + \frac{1}{c} = 0$, impossible.

\textbf{Subcase $a = 2$}: $\frac{1}{b} + \frac{1}{c} = \frac{1}{2}$ with $2 \leq b \leq c \leq 9$

Need: $\frac{1}{b} + \frac{1}{c} = \frac{1}{2}$ and $\frac{1}{b} + \frac{1}{c} > \frac{1}{2}$

This is a contradiction! So no solutions here.

\textbf{Subcase $a = 3$}: $\frac{1}{b} + \frac{1}{c} = \frac{2}{3}$ with $3 \leq b \leq c \leq 9$

Try $b = 3$: $\frac{1}{c} = \frac{2}{3} - \frac{1}{3} = \frac{1}{3}$, so $c = 3$.

Check ordering: $3 \leq 3 \leq 3 \leq 9$ \checkmark

Check triangle inequality: $\frac{1}{3} + \frac{1}{3} = \frac{2}{3} > \frac{1}{3}$ \checkmark

Solution: $(3, 3, 3)$ with $r = 1$ \checkmark

Try $b = 4$: $\frac{1}{c} = \frac{2}{3} - \frac{1}{4} = \frac{5}{12}$, so $c = \frac{12}{5}$ (not an integer).

For $b \geq 4$: $\frac{1}{b} \leq \frac{1}{4}$, so $\frac{1}{c} = \frac{2}{3} - \frac{1}{b} \geq \frac{2}{3} - \frac{1}{4} = \frac{5}{12}$

This means $c \leq \frac{12}{5} = 2.4 < b$, violating $b \leq c$.

\textbf{Conclusion for $r = 1$}: Only $(3, 3, 3)$ works.

\textbf{Case 2: $r = 2$}

We need: $\frac{1}{a} + \frac{1}{b} + \frac{1}{c} = \frac{1}{2}$ with $a \leq b \leq c \leq 9$

Since $\frac{1}{a} + \frac{1}{b} + \frac{1}{c} \leq \frac{3}{a} = \frac{1}{2}$:
$$a \leq 6$$

Also, for triangle inequality to not be violated immediately:
$$\frac{1}{b} + \frac{1}{c} > \frac{1}{a}$$

But from the formula:
$$\frac{1}{b} + \frac{1}{c} = \frac{1}{2} - \frac{1}{a}$$

So we need:
$$\frac{1}{2} - \frac{1}{a} > \frac{1}{a} \implies \frac{1}{2} > \frac{2}{a} \implies a > 4$$

Therefore: $5 \leq a \leq 6$.

\textbf{Subcase $a = 5$}: $\frac{1}{b} + \frac{1}{c} = \frac{1}{2} - \frac{1}{5} = \frac{3}{10}$ with $5 \leq b \leq c \leq 9$

Try $b = 5$: $\frac{1}{c} = \frac{3}{10} - \frac{1}{5} = \frac{1}{10}$, so $c = 10$.

But $c \leq 9$, so this doesn't work.

For $b \geq 5$: We need $c \geq 10$ (verified by algebra), violating $c \leq 9$.

\textbf{Subcase $a = 6$}: $\frac{1}{b} + \frac{1}{c} = \frac{1}{2} - \frac{1}{6} = \frac{1}{3}$ with $6 \leq b \leq c \leq 9$

Try $b = 6$: $\frac{1}{c} = \frac{1}{3} - \frac{1}{6} = \frac{1}{6}$, so $c = 6$.

Check ordering: $6 \leq 6 \leq 6 \leq 9$ \checkmark

Check triangle inequality: $\frac{1}{6} + \frac{1}{6} = \frac{1}{3} > \frac{1}{6}$ \checkmark

Solution: $(6, 6, 6)$ with $r = 2$ \checkmark

Try $b = 7, 8, 9$: Similar analysis shows $c$ exceeds bounds or gives non-integer.

\textbf{Conclusion for $r = 2$}: Only $(6, 6, 6)$ works.

\textbf{Case 3: $r = 3$}

We need: $\frac{1}{a} + \frac{1}{b} + \frac{1}{c} = \frac{1}{3}$ with $a \leq b \leq c \leq 9$

From $\frac{3}{a} \geq \frac{1}{3}$: $a \leq 9$.

For triangle inequality:
$$\frac{1}{b} + \frac{1}{c} = \frac{1}{3} - \frac{1}{a} > \frac{1}{a}$$
$$\frac{1}{3} > \frac{2}{a} \implies a > 6$$

Therefore: $7 \leq a \leq 9$.

By similar analysis to previous cases, only the symmetric solution works:

Try $a = 9$: $\frac{1}{b} + \frac{1}{c} = \frac{1}{3} - \frac{1}{9} = \frac{2}{9}$ with $9 \leq b \leq c \leq 9$

Only possibility: $b = c = 9$

Check: $\frac{1}{9} + \frac{1}{9} = \frac{2}{9}$ \checkmark

Solution: $(9, 9, 9)$ with $r = 3$ \checkmark

\textbf{Conclusion for $r = 3$}: Only $(9, 9, 9)$ works.

\textbf{Step 5: Final Count}

Valid triples: $(3, 3, 3)$, $(6, 6, 6)$, $(9, 9, 9)$

Total count: $\boxed{3}$

Answer: $\boxed{\textbf{(B) } 3}$

\end{solutionbox}

\begin{warning}
\textbf{Common Pitfalls}:
\begin{itemize}
    \item \textbf{Formula Derivation Error}: Carefully derive $\frac{1}{r} = \frac{1}{a} + \frac{1}{b} + \frac{1}{c}$ from area formulas
    \item \textbf{Forgetting Triangle Inequality}: Must check $\frac{1}{b} + \frac{1}{c} > \frac{1}{a}$ for each candidate
    \item \textbf{Wrong Bound on $r$}: Use $\frac{1}{r} \geq \frac{3}{c} \geq \frac{3}{9}$ to get $r \leq 3$
    \item \textbf{Missing Cases}: Systematically check $a = 1, 2, \ldots$ for each $r$
    \item \textbf{Arithmetic Errors}: Carefully compute reciprocal sums
    \item \textbf{Ordering Violation}: Ensure $a \leq b \leq c$ is maintained
    \item \textbf{Assuming Only Symmetric Cases}: Must verify that non-symmetric cases don't work
\end{itemize}
\end{warning}

\begin{protip}
\textbf{Time-Saving Strategies}:
\begin{itemize}
    \item \textbf{Pattern Recognition}: Notice $(3, 3, 3)$, $(6, 6, 6)$, $(9, 9, 9)$ pattern immediately
    \item \textbf{Bounds First}: Establish $r \leq 3$ before case analysis to limit work
    \item \textbf{Triangle Inequality Shortcut}: For $a = b = c$, inequality is automatic: $\frac{2}{a} > \frac{1}{a}$
    \item \textbf{Symmetric Guess}: Start with symmetric cases as they're most likely
    \item \textbf{Quick Verification}: For $a = b = c = 3k$, we get $\frac{1}{r} = \frac{3}{3k} = \frac{1}{k}$, so $r = k$
    \item \textbf{Answer Choice Leverage}: Answer (B) = 3 suggests looking for 3 solutions
\end{itemize}
\end{protip}

\begin{competitionbox}
\subsection*{7. Competition Strategy Integration}

\subsubsection*{A. Time Management}

\textbf{Time Budget}: 7-10 minutes for this P24 problem

\textbf{Allocation}:
\begin{itemize}
    \item Problem reading and understanding: 1 minute
    \item Formula derivation: 2-3 minutes
    \item Establish bounds on $r$: 1 minute
    \item Case analysis ($r = 1, 2, 3$): 3-4 minutes
    \item Verification: 1-2 minutes
\end{itemize}

\textbf{Decision Points}:
\begin{itemize}
    \item If derivation takes >4 minutes, try pattern recognition approach
    \item If case analysis stalls, guess symmetric solutions
    \item If time limited, verify answer (B) = 3 works with $(3,3,3), (6,6,6), (9,9,9)$
\end{itemize}

\subsubsection*{B. Prioritization Logic}

\textbf{Why Attempt This Problem}:
\begin{itemize}
    \item Clear mathematical structure (derive formula, case analysis)
    \item Small answer choices (2-6) suggest manageable work
    \item Pattern is recognizable (multiples of 3)
    \item Worth 6 points on AMC 12
    \item Formula derivation is instructive
\end{itemize}

\textbf{When to Skip}:
\begin{itemize}
    \item If uncomfortable deriving formulas from first principles
    \item If running very low on time (<10 minutes)
    \item If Problem 25 appears more tractable
    \item If weak with reciprocals and fractions
\end{itemize}

\subsubsection*{C. Risk Management}

\textbf{Confidence Indicators}:
\begin{itemize}
    \item Successfully derived formula (80\% confidence)
    \item Established bounds correctly (85\% confidence)
    \item Found pattern $(3k, 3k, 3k)$ (90\% confidence)
    \item Verified all three cases (95\% confidence)
    \item Checked no other solutions exist (98\% confidence)
\end{itemize}

\textbf{Fallback Strategies}:
\begin{itemize}
    \item If derivation unclear, guess symmetric pattern
    \item Test $(3,3,3), (6,6,6), (9,9,9)$ directly
    \item If stuck, answer (B) based on pattern recognition
\end{itemize}

\subsubsection*{D. Strategic Abandonment}

\textbf{Abandon If}:
\begin{itemize}
    \item Cannot derive formula after 4-5 minutes
    \item Case analysis becomes too complex (>15 subcases)
    \item Time exceeds 12 minutes without clear answer
    \item Getting lost in algebraic manipulations
\end{itemize}

\textbf{Before Abandoning}:
\begin{itemize}
    \item Try the pattern guess: $(3,3,3), (6,6,6), (9,9,9)$
    \item Verify these three work with the formula
    \item Answer (B) = 3 based on this
\end{itemize}

\end{competitionbox}

\section*{Summary}

\subsection*{Key Takeaways}

\begin{enumerate}
    \item \textbf{Formula Derivation}: The relationship $\frac{1}{r} = \frac{1}{a} + \frac{1}{b} + \frac{1}{c}$ is the key to the problem and can be derived from area formulas.
    
    \item \textbf{Bounds Analysis}: Using $c \leq 9$ gives $\frac{1}{r} \geq \frac{1}{3}$, limiting $r$ to at most 3. This dramatically reduces the search space.
    
    \item \textbf{Triangle Inequality for Reciprocals}: The condition $\frac{1}{b} + \frac{1}{c} > \frac{1}{a}$ ensures a non-degenerate triangle exists with the given altitudes.
    
    \item \textbf{Symmetric Solutions Only}: All valid solutions turn out to be symmetric: $(3,3,3)$, $(6,6,6)$, $(9,9,9)$.
    
    \item \textbf{Pattern Recognition}: The solutions are $(3k, 3k, 3k)$ for $k = 1, 2, 3$, which give inradii $r = k$.
    
    \item \textbf{Systematic Case Analysis}: Checking each value of $r$ exhaustively ensures no solutions are missed.
\end{enumerate}

\subsection*{Skills Practiced}

\begin{itemize}
    \item Deriving formulas from geometric principles
    \item Working with inradius and area formulas
    \item Reciprocal relationships and harmonic-like sums
    \item Establishing bounds from constraints
    \item Systematic case analysis
    \item Triangle inequality applications
    \item Integer solution finding (Diophantine-like)
    \item Counting with constraints
    \item Pattern recognition
\end{itemize}

\subsection*{Final Answer}

The number of valid ordered triples $(a, b, c)$ is:

\begin{center}
\fbox{\textbf{Answer: (B) 3}}
\end{center}

\textbf{The Three Solutions}:
\begin{itemize}
    \item $(a, b, c) = (3, 3, 3)$ with inradius $r = 1$
    \item $(a, b, c) = (6, 6, 6)$ with inradius $r = 2$
    \item $(a, b, c) = (9, 9, 9)$ with inradius $r = 3$
\end{itemize}

All are verified to satisfy:
\begin{itemize}
    \item Formula: $\frac{1}{r} = \frac{1}{a} + \frac{1}{b} + \frac{1}{c}$ \checkmark
    \item Bounds: $a \leq b \leq c \leq 9$ \checkmark
    \item Triangle inequality: $\frac{1}{b} + \frac{1}{c} > \frac{1}{a}$ \checkmark
\end{itemize}

\end{document}